% GENERATED FILE. Edit canonical lessons, then run npm run generate:books.
% canonical-source-hash: fnv1a64:8cf375e39a133731
% canonical-lessons: IT-C257-metallo, IT-C257-mattone, IT-C257-sacchetto, IT-C257-filo, IT-C257-strumento

\chapter{Metal, Brick, Small bag, Thread, Tool}
\label{ch:it-metal-brick-small-bag-thread-tool}
\hlchaptermodality{257}

\begin{chapteropening}
\textbf{By the end of this chapter:} \emph{I can say metal, a brick, a small bag, a thread and a tool.}

Complete the last lesson of chapter 257: I can say metal, a brick, a small bag, a thread and a tool.
\end{chapteropening}

\section[il metallo]{il metallo — metal}

\label{lesson:IT-C257-metallo}



\begin{quote}
Before the new one: say the Italian for a rainbow, then the Italian for the climate.
\end{quote}

\subsection*{You'll want to know: il metallo}
\textbf{il metallo} — \textquotedblleft{}metal\textquotedblright{}.

\begin{cousinweb}
\textbf{metallo} $\leftarrow$ Latin \emph{metallum}, \textquotedblleft{}a mine, a metal\textquotedblright{}, from Greek \emph{metallon} — a learned borrowing. (EtymDB 2.1 records the Latin and Greek words.)
\end{cousinweb}

\begin{grammarlens}[title={What you've built}]
One more word for this chapter.
\end{grammarlens}

\subsection*{Your turn}
\begin{itemize}
\raggedright
  \item \emph{Say it:} \emph{il metallo}
  \item \emph{Say it:} \emph{il metallo}, once more
  \item \emph{Say it:} say \emph{il metallo}
  \item \emph{Recall:} say the Italian for a rainbow, then the Italian for the climate, then say \emph{il metallo} again
\end{itemize}

\begin{tcolorbox}[breakable,colback=teal!4,colframe=teal!35!black,title={Before you move on}]
What does il metallo mean? (\textquotedblleft{}Metal\textquotedblright{}.) Say it once more.
\end{tcolorbox}
\section[il mattone]{il mattone — a brick}

\label{lesson:IT-C257-mattone}



\begin{quote}
Before the new one: say the Italian for the climate, then the Italian for metal.
\end{quote}

\subsection*{You'll want to know: il mattone}
\textbf{il mattone} — \textquotedblleft{}a brick\textquotedblright{}.

\begin{cousinweb}
\textbf{mattone} — its origin is uncertain. EtymDB 2.1 records the word but not where it came from.
\end{cousinweb}

\begin{grammarlens}[title={What you've built}]
One more word for this chapter.
\end{grammarlens}

\subsection*{Your turn}
\begin{itemize}
\raggedright
  \item \emph{Say it:} \emph{il mattone}
  \item \emph{Say it:} \emph{il mattone}, once more
  \item \emph{Say it:} say \emph{il mattone}
  \item \emph{Recall:} say the Italian for the climate, then the Italian for metal, then say \emph{il mattone} again
\end{itemize}

\begin{tcolorbox}[breakable,colback=teal!4,colframe=teal!35!black,title={Before you move on}]
What does il mattone mean? (\textquotedblleft{}A brick\textquotedblright{}.) Say it once more.
\end{tcolorbox}
\section[il sacchetto]{il sacchetto — a small bag}

\label{lesson:IT-C257-sacchetto}



\begin{quote}
Before the new one: say the Italian for metal, then the Italian for a brick.
\end{quote}

\subsection*{You'll want to know: il sacchetto}
\textbf{il sacchetto} — \textquotedblleft{}a small bag\textquotedblright{}.

\begin{cousinweb}
\textbf{sacchetto} — built from \emph{sacco}, \textquotedblleft{}a sack\textquotedblright{}, with \emph{-etto}, \textquotedblleft{}little\textquotedblright{}. \emph{sacco} $\leftarrow$ Latin \emph{saccus}, from Greek \emph{sakkos}, \textquotedblleft{}a bag of coarse cloth\textquotedblright{}. (EtymDB 2.1)
\end{cousinweb}

\begin{grammarlens}[title={What you've built}]
One more word for this chapter.
\end{grammarlens}

\subsection*{Your turn}
\begin{itemize}
\raggedright
  \item \emph{Say it:} \emph{il sacchetto}
  \item \emph{Say it:} \emph{il sacchetto}, once more
  \item \emph{Say it:} say \emph{il sacchetto}
  \item \emph{Recall:} say the Italian for metal, then the Italian for a brick, then say \emph{il sacchetto} again
\end{itemize}

\begin{tcolorbox}[breakable,colback=teal!4,colframe=teal!35!black,title={Before you move on}]
What does il sacchetto mean? (\textquotedblleft{}A small bag\textquotedblright{}.) Say it once more.
\end{tcolorbox}
\section[il filo]{il filo — a thread, a wire}

\label{lesson:IT-C257-filo}



\begin{quote}
Before the new one: say the Italian for a brick, then the Italian for a small bag.
\end{quote}

\subsection*{You'll want to know: il filo}
\textbf{il filo} — \textquotedblleft{}a thread, a wire\textquotedblright{}.

\begin{cousinweb}
\textbf{filo} $\leftarrow$ Latin \emph{filum}, \textquotedblleft{}a thread\textquotedblright{} — inherited. (EtymDB 2.1)
\end{cousinweb}

\begin{grammarlens}[title={What you've built}]
One more word for this chapter.
\end{grammarlens}

\subsection*{Your turn}
\begin{itemize}
\raggedright
  \item \emph{Say it:} \emph{il filo}
  \item \emph{Say it:} \emph{il filo}, once more
  \item \emph{Say it:} say \emph{il filo}
  \item \emph{Recall:} say the Italian for a brick, then the Italian for a small bag, then say \emph{il filo} again
\end{itemize}

\begin{tcolorbox}[breakable,colback=teal!4,colframe=teal!35!black,title={Before you move on}]
What does il filo mean? (\textquotedblleft{}A thread, a wire\textquotedblright{}.) Say it once more.
\end{tcolorbox}
\section[lo strumento]{lo strumento — a tool, an instrument}

\label{lesson:IT-C257-strumento}



\begin{quote}
Before the new one: say the Italian for a small bag, then the Italian for a thread.
\end{quote}

\subsection*{You'll want to know: lo strumento}
\textbf{lo strumento} — \textquotedblleft{}a tool, an instrument\textquotedblright{}.

\begin{cousinweb}
\textbf{strumento} $\leftarrow$ Latin \emph{instrumentum}, \textquotedblleft{}a tool\textquotedblright{}, its first syllable worn away — inherited. EtymDB 2.1 has no row for \emph{strumento}; it derives the older Italian \emph{istrumento} from \emph{instrumentum}.
\end{cousinweb}

\begin{grammarlens}[title={What you've built}]
One more word for this chapter.
\end{grammarlens}

\subsection*{Your turn}
\begin{itemize}
\raggedright
  \item \emph{Say it:} \emph{lo strumento}
  \item \emph{Say it:} \emph{lo strumento}, once more
  \item \emph{Say it:} say \emph{lo strumento}
  \item \emph{Recall:} say the Italian for a small bag, then the Italian for a thread, then say \emph{lo strumento} again
\end{itemize}

\begin{tcolorbox}[breakable,colback=teal!4,colframe=teal!35!black,title={Before you move on}]
What does lo strumento mean? (\textquotedblleft{}A tool, an instrument\textquotedblright{}.) Say it once more.
\end{tcolorbox}
